\documentclass[12pt,a4paper]{article}
\usepackage{mathwizards-ingles}
\title{%
  \vspace{-1cm}
  {\Huge\bfseries Programa de Inmersión en Inglés}\\[0.3cm]
  {\Large De Cero Absoluto a Usuario Independiente}\\[0.2cm]
  {\large\color{grissuave} Nivel A0 $\rightarrow$ B1 · Marco Común Europeo de Referencia}\\[0.5cm]
  \rule{\textwidth}{1.5pt}
}
\author{}
\date{}

\begin{document}
\maketitle
\thispagestyle{empty}

\vspace{-0.5cm}

% ═══════════════════════════════════════════════════════════════════════════════
\section*{Visión General del Curso}
% ═══════════════════════════════════════════════════════════════════════════════

\begin{cajafilosofia}[Filosofía Pedagógica]
Este curso se fundamenta en dos pilares de la lingüística aplicada contemporánea:

\begin{enumerate}[leftmargin=1.5cm]
  \item \textbf{Hipótesis del Input Comprensible (CI)} de Stephen Krashen: la adquisición ocurre cuando el estudiante comprende mensajes en la lengua meta que contienen estructuras ligeramente por encima de su nivel actual (\textit{i+1}).

  \item \textbf{Sistemas de Repetición Espaciada (SRS)}: optimización del procesamiento profundo mediante herramientas como Anki, que consolidan la memoria a largo plazo distribuyendo los repasos según la curva del olvido.
\end{enumerate}

\medskip
\textbf{Modalidad:} B-Learning (Blended Learning) virtual.\\
\textbf{Duración:} 12 semanas (3 meses).\\
\textbf{Frecuencia:} 4--5 sesiones semanales (3 sincrónicas + trabajo asincrónico diario).\\
\textbf{Compromiso diario estimado:} 1.5 a 2 horas.
\end{cajafilosofia}

% ═══════════════════════════════════════════════════════════════════════════════
\section*{Roles y Responsabilidades}
% ═══════════════════════════════════════════════════════════════════════════════

\begin{multicols}{2}

\begin{cajames}[Rol del Profesor/Moderador]{acento}
\begin{itemize}[leftmargin=0.8cm]
  \item Reducir el filtro afectivo
  \item Seleccionar y proveer material adecuado al nivel
  \item Facilitar sesiones sincrónicas con TPR, ALG, lectura guiada
  \item Validar oraciones minadas (regla \textit{i+1})
  \item Modelar pronunciación y reformular errores (\textit{recasting})
  \item Monitorear el bienestar del grupo
\end{itemize}
\end{cajames}

\columnbreak

\begin{cajames}[Rol del Estudiante]{acento}
\begin{itemize}[leftmargin=0.8cm]
  \item Asistir a sesiones sincrónicas
  \item Completar inmersión asincrónica diaria
  \item Revisar tarjetas Anki diariamente (5--10 nuevas/día)
  \item Participar en foros grupales (Discord/Slack)
  \item Compartir oraciones minadas para validación
  \item Tolerar la ambigüedad sin recurrir a traducción
\end{itemize}
\end{cajames}

\end{multicols}

% ═══════════════════════════════════════════════════════════════════════════════
\section*{Arquitectura del Programa: 3 Meses}
% ═══════════════════════════════════════════════════════════════════════════════

% ── MES 1 ────────────────────────────────────────────────────────────────────
\begin{cajames}[\textcolor{white}{Mes 1: El Periodo Silencioso y Adquisición Fonológica (A0 $\rightarrow$ A1)}]{mes1}

\textbf{Semanas 1--4} · 12 sesiones sincrónicas

\medskip
\begin{tabularx}{\textwidth}{@{} l X @{}}
\toprule
\textbf{Componente} & \textbf{Descripción} \\
\midrule
Calentamiento Visual y TPR & El profesor actúa, gesticula y usa material visual. Respuesta Física Total. Los estudiantes responden sin forzar el habla. \\[4pt]
Inmersión Narrativa & Videos de ``Súper Principiantes'' y videoblogs \textit{Slice of Life}. Sin subtítulos en español. \\[4pt]
Estudio Deliberado & Mazo Anki prefabricado compartido por el profesor. Vocabulario de alta frecuencia. \\
\bottomrule
\end{tabularx}

\medskip
\textbf{Gramática emergente:} Verbos de acción (TPR), sustantivos concretos, saludos, presentaciones.\\
\textbf{Habilidad principal:} Comprensión auditiva con soporte visual.\\
\textbf{Producción esperada:} Nula o mínima (reacciones no verbales).
\end{cajames}

\vspace{0.3cm}

% ── MES 2 ────────────────────────────────────────────────────────────────────
\begin{cajames}[\textcolor{white}{Mes 2: Minería Colaborativa y Lectura Asistida (A1 $\rightarrow$ A2)}]{mes2}

\textbf{Semanas 5--8} · 12 sesiones sincrónicas

\medskip
\begin{tabularx}{\textwidth}{@{} l X @{}}
\toprule
\textbf{Componente} & \textbf{Descripción} \\
\midrule
Círculos de Lectura Virtuales & Lectura guiada de \textit{Graded Readers} (A1/A2). Inferencia colaborativa. Modelo fonético en vivo. \\[4pt]
Inmersión Visual Narrativa & Transición a acentos diversos. Videos nivel A2. \\[4pt]
Minería de Oraciones Colaborativa & Extracción de oraciones (\textit{i+1}). Validación grupal moderada por el profesor. \\[4pt]
Revisión Espaciada (SRS) & Tarjetas personalizadas + curadas grupalmente en Anki. \\
\bottomrule
\end{tabularx}

\medskip
\textbf{Gramática emergente:} Present Simple, Past Simple, Present Continuous, comparativos.\\
\textbf{Habilidad principal:} Lectura asistida + comprensión auditiva.\\
\textbf{Producción esperada:} Frases cortas, respuestas a preguntas simples.
\end{cajames}

\vspace{0.3cm}

% ── MES 3 ────────────────────────────────────────────────────────────────────
\begin{cajames}[\textcolor{white}{Mes 3: Consolidación Auditiva y Clubes de Debate (A2 $\rightarrow$ B1)}]{mes3}

\textbf{Semanas 9--12} · 12 sesiones sincrónicas

\medskip
\begin{tabularx}{\textwidth}{@{} l X @{}}
\toprule
\textbf{Componente} & \textbf{Descripción} \\
\midrule
Clubes de Escucha y Debate & Discusión de podcasts. Transición hacia el habla. El profesor reformula suavemente (\textit{recasting}). \\[4pt]
Inmersión Auditiva y Lectura & Consumo intensivo de podcasts y textos B1. Tolerancia a la ambigüedad. \\[4pt]
Minería Avanzada & Enfoque en \textit{collocations}, modismos e idioms. Peer-review de estructuras complejas. \\
\bottomrule
\end{tabularx}

\medskip
\textbf{Gramática emergente:} Present Perfect, Past Continuous, reported speech, conditionals, passive voice.\\
\textbf{Habilidad principal:} Comprensión auditiva pura + producción oral incipiente.\\
\textbf{Producción esperada:} Participación en debates, opiniones estructuradas.
\end{cajames}

% ═══════════════════════════════════════════════════════════════════════════════
\section*{Distribución Semanal Tipo}
% ═══════════════════════════════════════════════════════════════════════════════

\begin{cajanota}[Ejemplo: Semana típica del Mes 2]
\begin{tabularx}{\textwidth}{@{} c l X @{}}
\toprule
\textbf{Día} & \textbf{Modalidad} & \textbf{Actividad} \\
\midrule
Lunes    & \textcolor{mes2}{\textbf{Sincrónica}} & Círculo de Lectura Virtual \\
Martes   & Asincrónica & Inmersión visual (30 min) + Anki (15 min) \\
Miércoles & \textcolor{mes2}{\textbf{Sincrónica}} & Taller de Minería Colaborativa \\
Jueves   & Asincrónica & Lectura extensiva (30 min) + Anki (15 min) \\
Viernes  & \textcolor{mes2}{\textbf{Sincrónica}} & Sesión de Crowdsourcing de tarjetas \\
Sábado   & Asincrónica & Inmersión libre (video/lectura) + Anki (15 min) \\
Domingo  & Descanso & Solo repaso Anki de mantenimiento (10 min) \\
\bottomrule
\end{tabularx}
\end{cajanota}

% ═══════════════════════════════════════════════════════════════════════════════
\section*{Principios Innegociables}
% ═══════════════════════════════════════════════════════════════════════════════

\begin{enumerate}[leftmargin=1.5cm, label=\textbf{\arabic*.}]
  \item \textbf{Cero subtítulos en español.} Tolerancia absoluta a la ambigüedad.
  \item \textbf{Máximo 5--10 tarjetas nuevas por día en Anki.} Prevención de burnout.
  \item \textbf{Regla del \textit{i+1}.} Cada oración minada debe tener \textbf{una sola} incógnita.
  \item \textbf{No traducción sistemática.} Definiciones en inglés a partir del final del Mes 2.
  \item \textbf{Algoritmo FSRS en Anki.} Erradicación de ``sanguijuelas'' (tarjetas que fallan $>5$ veces).
  \item \textbf{Participación grupal activa.} La comunidad reduce el filtro afectivo y mantiene la motivación.
\end{enumerate}

% ═══════════════════════════════════════════════════════════════════════════════
\section*{Recursos del Curso}
% ═══════════════════════════════════════════════════════════════════════════════

\begin{cajanota}[Canales y Plataformas Recomendados]
\begin{multicols}{2}
\textbf{Videos CI (Beginners):}
\begin{itemize}[leftmargin=0.8cm, nosep]
  \item Comprehensible English
  \item English Fiesta
  \item Fingtam English
  \item EnglishSponge (A2)
\end{itemize}

\textbf{Lectura:}
\begin{itemize}[leftmargin=0.8cm, nosep]
  \item English e-Reader (.net)
  \item Oxford Graded Readers
  \item Pearson English Readers
\end{itemize}

\columnbreak

\textbf{Podcasts:}
\begin{itemize}[leftmargin=0.8cm, nosep]
  \item Easy Stories in English
  \item Learn English Podcast
  \item VOA Learning English
\end{itemize}

\textbf{Herramientas:}
\begin{itemize}[leftmargin=0.8cm, nosep]
  \item Anki + FSRS + AnkiConnect
  \item Language Reactor (ext.)
  \item Yomitan (ext.)
  \item Discord/Slack (comunidad)
\end{itemize}
\end{multicols}
\end{cajanota}

\vfill
\begin{center}
\rule{0.5\textwidth}{0.5pt}\\[0.3cm]
{\small\color{grissuave} Diseño curricular basado en la metodología Refold · Hipótesis del Input Comprensible (Krashen) · MCER}
\end{center}

\end{document}
